\begin{center}
  \large{\textbf{Abstract}}
\end{center}
\selectlanguage{english}
\begingroup
\setstretch{1.2}
\noindent What do scientists and engineers do when they develop artificial intelligence? Driven by this question, I accompanied researchers over four years of ethnographic work at the Artificial Intelligence Center of the University of São Paulo, one of the Applied Research Centers of the São Paulo Research Foundation. The center was created in 2020 through an agreement between the foundation, the university, and the International Business Machines Corporation. When I began the research, in 2022, the partnership appeared stable; when I finished it, in 2025, the corporation had ended the agreement. The extended fieldwork period made it possible to describe the arrangement of technoscientific practices within the sociotechnical network that linked the public university and the transnational corporation in artificial intelligence research in Brazil, from its foundation to its dissolution. By following Fábio Cozman and Cláudio Pinhanez, the center's directors, the research traversed the history of the relationship between the center and the corporation and the corporation's technoscientific practices over time. By following Marcelo Finger, coordinator of the project Early Detection System for Respiratory Failure based on Audio Analysis, I traced the chain of translations that converted speech recordings of patients hospitalized with Covid-19 into visual representations of sound, processed by neural networks for the detection of respiratory failure through voice analysis. This ethnographic gesture of accompanying actors wherever they went led the thesis through multiple sites, and it was this movement that led me to choose actor-network theory as the theoretical and methodological orientation capable of accounting for the field. Along the way, a lexicometric analysis of works by Bruno Latour situated the conceptual vocabulary I mobilize, and a bibliometric mapping of studies on artificial intelligence in the Brazilian humanities situated this thesis within the field in which it intervenes. Throughout the research, carried out not only \textit{about} but \textit{with} generative artificial intelligence, the production of ethnographic knowledge mobilized a set of technoscientific inscription practices: interactive diagrams, computational analyses, artificial intelligence, charts, graphs, programming languages, and the construction of a website for an interactive reading of the thesis. These technoscientific inscriptions constitute the concept of technographies that I propose, and they designate a mode of research that inhabits the tension between the technique of inscriptions and the sensible reality of human bodies. By following actors in an artificial intelligence research center, this thesis contributes to science and technology studies by describing, within a single trajectory, the institutional and corporate negotiations that sustain artificial intelligence research and the situated technoscientific practices of system development, and by showing how these two planes articulate.

\noindent \textbf{Keywords}: Artificial Intelligence; Ethnography; Technographies; Technoscientific Practices; Science and Technology Studies.
\endgroup
\selectlanguage{brazilian}